\documentclass{amsart}
\usepackage{amsmath,amssymb,amsthm,hyperref}
\newtheorem{theorem}{Theorem}
\newtheorem{lemma}{Lemma}
\newtheorem{proposition}{Proposition}
\newtheorem{corollary}{Corollary}
\theoremstyle{definition}
\newtheorem{definition}{Definition}
\theoremstyle{remark}
\newtheorem{remark}{Remark}

\begin{document}

\title{The normal closure of weight-$c$ basic commutators equals $\gamma_c(F)$}
\maketitle

\section{Statement and conventions}

Let $F$ be a free group of finite rank $r$ with fixed free generating set
$X=\{x_1,\dots,x_r\}$. We write $X^{\pm}=X\cup X^{-1}$. We use the commutator
convention
\[
[a,b]=a^{-1}b^{-1}ab ,\qquad a^{b}=b^{-1}ab .
\]
The lower central series is $\gamma_1(F)=F$ and $\gamma_{k+1}(F)=[\gamma_k(F),F]$
for $k\ge 1$; for subgroups $H,K\le F$,
\[
[H,K]=\big\langle\, [h,k] : h\in H,\ k\in K \,\big\rangle .
\]
Each $\gamma_k(F)$ is fully invariant, hence normal, and
$\gamma_{k+1}(F)\le\gamma_k(F)$.

Basic commutators in $X$ and their weights are defined by the Hall ordering as in
the problem statement. We fix once and for all a total order
$c_1<c_2<c_3<\cdots$ of the (countable) set of all basic commutators, refining
weight (smaller weight precedes larger), with $c_1=x_1,\dots,c_r=x_r$ the basic
commutators of weight $1$. For $c\ge 1$ let
\[
B_c=\{\,b : b \text{ is a basic commutator of weight exactly } c\,\},
\qquad
N_c=\langle B_c\rangle^{F},
\]
where $\langle S\rangle^{F}$ denotes the normal closure of $S$ in $F$. We write
$B_{\ge c}=\bigcup_{m\ge c}B_m$.

\begin{theorem}\label{thm:main}
For every integer $c\ge 1$ one has $\gamma_c(F)=N_c$. Equivalently, the answer to
the problem is \emph{yes}.
\end{theorem}

\paragraph{Plan and honest accounting of inputs.}
The inclusion $N_c\subseteq\gamma_c(F)$ is elementary (\S\ref{sec:easy}). The
reverse inclusion $\gamma_c(F)\subseteq N_c$ is the substance of the theorem. We
prove it in three steps.

\begin{enumerate}
\item[(I)] An elementary \emph{chaining reduction} (\S\ref{sec:chain},
Proposition~\ref{prop:chain}) showing that the entire theorem
($\gamma_c=N_c$ for all $c$) is equivalent to the single-level statement
\begin{equation}\label{eq:S}
\text{\emph{(S)}}\qquad B_w\subseteq N_{w-1}\quad\text{for every }w\ge 2 .
\end{equation}
\item[(II)] A rigorous proof that $\gamma_c(F)\subseteq N_c$ \emph{follows from}
(S), using P.~Hall's collection theorem (\S\ref{sec:collect}): every element of
$\gamma_c(F)$ is a \emph{finite} ordered product of powers of basic commutators of
weight $\ge c$, and (S) (via chaining) puts each such factor in $N_c$. The
\emph{finiteness} of the collected form is what makes this step exact; it removes
the ``infinite product / closure'' obstruction that defeats every purely graded
argument (\S\ref{sec:why}).
\item[(III)] The statement (S) itself. We isolate it (\S\ref{sec:S}) as the one
genuinely non-elementary fact, prove that it cannot be obtained from the graded
basis theorem plus elementary commutator calculus (\S\ref{sec:why}: the
obstruction is residual nilpotence of $F/N_c$, which is not a graded invariant),
and cite for it the standard collection-process proof of the Hall basis theorem.
We are explicit that (S) is the substantive classical input and we do not disguise
it as a consequence of weaker primitives.
\end{enumerate}

\section{The classical inputs, stated precisely}\label{sec:input}

We use the Hall basis theorem in two standard forms. Neither of the two cited
external statements is the conclusion of Theorem~\ref{thm:main}; the conclusion
is \emph{deduced} from them in \S\S\ref{sec:chain}--\ref{sec:collect}, with the
single substantive use isolated in \S\ref{sec:S}.

\begin{proposition}[Hall basis theorem, graded form]\label{prop:hall}
For each $k\ge 1$ the basic commutators of weight $\ge k$ generate $\gamma_k(F)$
\emph{as a subgroup}, and $\gamma_k(F)/\gamma_{k+1}(F)$ is free abelian with basis
the images of the basic commutators of weight exactly $k$. In particular every
basic commutator of weight $k$ lies in $\gamma_k(F)$, and
\begin{equation}\label{eq:hall-decomp}
\gamma_k(F)=\langle B_k\rangle\cdot\gamma_{k+1}(F).
\end{equation}
\end{proposition}

\noindent(M.~Hall, \emph{The Theory of Groups}, Thm.~11.2.4;
Magnus--Karrass--Solitar, \emph{Combinatorial Group Theory}, Thm.~5.13A.) This is
the \emph{graded} statement.

\begin{theorem}[Hall collection theorem; finite collected normal form]\label{thm:collect}
There is a terminating rewriting procedure (the collecting process) with the
following property. Let $g\in F$ be represented by a word of length $L$ in
$X^{\pm}$. Then $g$ can be written as a \emph{finite} ordered product
\begin{equation}\label{eq:collected}
g=c_{i_1}^{\,a_1}c_{i_2}^{\,a_2}\cdots c_{i_t}^{\,a_t},
\qquad i_1<i_2<\cdots<i_t,\ a_j\in\mathbb Z,
\end{equation}
of powers of basic commutators, in which every basic commutator $c_{i_j}$ that
occurs has weight at most $L$; this is an \emph{identity in $F$} (not merely modulo
a term of the lower central series). Moreover, for each fixed $m$, the truncation
modulo $\gamma_m(F)$ of \eqref{eq:collected}, namely
$\prod_{\operatorname{wt}(c_{i_j})<m}\bar c_{i_j}^{\,a_j}$, is the
\emph{unique} expression of $\bar g$ in $F/\gamma_m(F)$ as an ordered product of
powers of (the images of) the basic commutators of weight $<m$; in particular the
exponents $a_j$ of the basic commutators of weight $<m$ are uniquely determined by
$g$.
\end{theorem}

\noindent(M.~Hall, \emph{The Theory of Groups}, Thm.~11.1.4 and \S11.2;
Magnus--Karrass--Solitar, \S5.3 and Thm.~5.13A.) The weight bound
``$\operatorname{wt}(c_{i_j})\le L$'' is the conservation of total weight under the
collecting moves: each move replaces an adjacent pair by a commutator whose weight
is the sum of the two weights, and the total weight of a word equals its length;
hence the process produces only finitely many basic commutators, all of weight
$\le L$, and \eqref{eq:collected} is a genuine finite identity in $F$. The
uniqueness clause is the freeness of $F/\gamma_m(F)$ on the basic commutators of
weight $<m$ (Proposition~\ref{prop:hall} layer by layer).

\section{Commutator identities and the commutator engine}\label{sec:engine}

We record the standard commutator identities, valid in any group with
$[a,b]=a^{-1}b^{-1}ab$:
\begin{align}
[ab,c]&=[a,c]^{b}\,[b,c], \label{eq:id1}\\
[a,bc]&=[a,c]\,[a,b]^{c}, \label{eq:id2}\\
[a,b^{-1}]&=\big([a,b]^{-1}\big)^{b^{-1}}, \label{eq:id3}\\
[a^{-1},b]&=\big([a,b]^{-1}\big)^{a^{-1}}, \label{eq:id4}\\
[a^{f},b]&=[a,\,b^{\,f^{-1}}]^{\,f}, \label{eq:id5}\\
[a,b]&=[b,a]^{-1}. \label{eq:id6}
\end{align}
Each is verified by direct expansion. (For \eqref{eq:id4}: put $a^{-1}$ for $a$ in
\eqref{eq:id1} to get $1=[a^{-1},c]^{a}[a,c]$. For \eqref{eq:id5}:
$[a^f,b]=f^{-1}a^{-1}f\,b^{-1}\,f^{-1}af\,b=[a,b^{f^{-1}}]^{f}$.)

\begin{lemma}[Commutator engine]\label{lem:engine}
Let $S\subseteq F$ and $H=\langle S\rangle^{F}$. Then $H$ is normal and
\begin{equation}\label{eq:engine}
[H,F]=\big\langle\,[s,x] : s\in S,\ x\in X^{\pm}\,\big\rangle^{F}.
\end{equation}
\end{lemma}
\begin{proof}
$H$ is normal by construction. Write
$M=\langle\,[s,x]:s\in S,\,x\in X^{\pm}\,\rangle^{F}$. Each $[s,x]\in[H,F]$, which is
normal; hence $M\subseteq[H,F]$.

For the reverse inclusion, $[H,F]$ is generated by the $[h,y]$ with $h\in H$,
$y\in F$; as $M$ is normal it suffices to prove $[h,y]\in M$ for all such.

\emph{Step 1: $[s,y]\in M$ for $s\in S$, $y\in F$.} Induct on the word length
$\ell$ of $y$ in $X^{\pm}$. For $\ell=0$, $[s,1]=1$. For $\ell=1$: $y=x\in X$ gives
$[s,x]\in M$; $y=x^{-1}$ gives $[s,x^{-1}]=([s,x]^{-1})^{x^{-1}}\in M$ by
\eqref{eq:id3} and normality. For $\ell\ge2$ write $y=y'z$, $z\in X^{\pm}$; by
\eqref{eq:id2}, $[s,y]=[s,z]\,[s,y']^{z}\in M$ by induction and normality.

\emph{Step 2: $[s^{f},x]\in M$ for $s\in S$, $f\in F$, $x\in X^{\pm}$.} By
\eqref{eq:id5}, $[s^{f},x]=[s,x^{f^{-1}}]^{f}\in M$ by Step~1 and normality.

\emph{Step 3: $[t,y]\in M$ for $t\in\{(s^{f})^{\pm1}\}$, $y\in F$.} For $t=s^f$,
Step~2 plus the word induction of Step~1 (verbatim with $s^f$ for $s$) gives
$[s^f,y]\in M$. For $t=(s^f)^{-1}$, \eqref{eq:id4} gives
$[(s^f)^{-1},x]=([s^f,x]^{-1})^{(s^f)^{-1}}\in M$, and the same word induction gives
$[(s^f)^{-1},y]\in M$.

\emph{Step 4: $[h,y]\in M$ for $h\in H$.} Write $h=t_1\cdots t_n$ with each
$t_i\in\{(s_i^{f_i})^{\pm1}\}$ and induct on $n$ via \eqref{eq:id1}:
$[h,y]=[t_1,y]^{h'}[h',y]$, $h'=t_2\cdots t_n$; both factors lie in $M$ by Step~3,
the inductive hypothesis, and normality.
\end{proof}

\begin{lemma}\label{lem:caseA}
If $u\in N_d$ (any $d$) and $v\in F$, then $[u,v]\in N_d$ and $[v,u]\in N_d$.
\end{lemma}
\begin{proof}
$N_d$ is normal, so $u^{v}\in N_d$ and $[u,v]=u^{-1}u^{v}\in N_d$; also
$[v,u]=(u^{-1})^{v}u\in N_d$.
\end{proof}

\section{The easy inclusion}\label{sec:easy}

\begin{lemma}\label{lem:easy}
For every $c\ge1$, $N_c\subseteq\gamma_c(F)$.
\end{lemma}
\begin{proof}
By Proposition~\ref{prop:hall} every $b\in B_c$ lies in the normal subgroup
$\gamma_c(F)$; hence $N_c=\langle B_c\rangle^{F}\subseteq\gamma_c(F)$.
\end{proof}

\section{Reducing the theorem to the single-level statement \eqref{eq:S}}
\label{sec:chain}

\begin{lemma}[Monotonicity from \eqref{eq:S}]\label{lem:mono}
If $B_w\subseteq N_{w-1}$ for some $w\ge2$, then $N_w\subseteq N_{w-1}$.
\end{lemma}
\begin{proof}
$N_{w-1}$ is a normal subgroup containing $B_w$, hence contains $N_w=\langle
B_w\rangle^{F}$.
\end{proof}

\begin{lemma}[$B_{\ge c}\subseteq N_c$ from \eqref{eq:S}]\label{lem:Bgec}
Assume \eqref{eq:S}. Then for every $c\ge1$ and every $w\ge c$, $B_w\subseteq N_c$;
hence $B_{\ge c}\subseteq N_c$.
\end{lemma}
\begin{proof}
Fix $c$ and induct on $w\ge c$. For $w=c$, $B_c\subseteq N_c$ by definition. For
$w>c$: by \eqref{eq:S} and Lemma~\ref{lem:mono}, $N_w\subseteq N_{w-1}$, so
$B_w\subseteq N_w\subseteq N_{w-1}$. By the inductive hypothesis (excess
$w-1-c<w-c$, $w-1\ge c$), $B_{w-1}\subseteq N_c$, hence
$N_{w-1}=\langle B_{w-1}\rangle^{F}\subseteq N_c$. Therefore $B_w\subseteq
N_{w-1}\subseteq N_c$.
\end{proof}

\begin{proposition}[Chaining reduction]\label{prop:chain}
The following are equivalent:
\begin{enumerate}
\item[\textup{(a)}] $\gamma_c(F)=N_c$ for every $c\ge1$ \textup{(Theorem~\ref{thm:main})};
\item[\textup{(b)}] $B_{\ge c}\subseteq N_c$ for every $c\ge1$;
\item[\textup{(c)}] the single-level statement \eqref{eq:S}: $B_w\subseteq N_{w-1}$ for every $w\ge2$.
\end{enumerate}
\end{proposition}
\begin{proof}
(a)$\Rightarrow$(b): if $\gamma_c(F)=N_c$ then each $b\in B_w$ ($w\ge c$) lies in
$\gamma_w(F)\subseteq\gamma_c(F)=N_c$ (Proposition~\ref{prop:hall}).

(b)$\Rightarrow$(c): take $c=w-1$; then $B_w\subseteq B_{\ge w-1}\subseteq N_{w-1}$.

(c)$\Rightarrow$(b) is Lemma~\ref{lem:Bgec}.

(b)$\Rightarrow$(a): For the inclusion $\gamma_c(F)\subseteq N_c$ we use the
finite collected form; this is carried out in
Proposition~\ref{prop:gamma-in-N} below. Combined with Lemma~\ref{lem:easy} this
gives $\gamma_c(F)=N_c$.
\end{proof}

\noindent To avoid any appearance of circularity we make the dependency explicit:
the implication (b)$\Rightarrow$(a) of Proposition~\ref{prop:chain} is
\emph{completed} by Proposition~\ref{prop:gamma-in-N} of the next section, which
uses (b) (equivalently (c)) but nothing about (a).

\section{From \eqref{eq:S} to $\gamma_c(F)\subseteq N_c$ via finite collected forms}
\label{sec:collect}

\begin{lemma}[Elements of $\gamma_c$ are finite products of high-weight basic commutators]\label{lem:finite-form}
Let $g\in\gamma_c(F)$. Then $g$ is a finite ordered product
\begin{equation}\label{eq:gform}
g=c_{i_1}^{\,a_1}\cdots c_{i_t}^{\,a_t}
\end{equation}
of powers of basic commutators, all of weight $\ge c$.
\end{lemma}
\begin{proof}
By Theorem~\ref{thm:collect}, $g$ has a finite collected form
\eqref{eq:collected}: $g=\prod_{j=1}^{t}c_{i_j}^{\,a_j}$ with all
$\operatorname{wt}(c_{i_j})\le L$ (where $L$ is the length of a word representing
$g$), the product taken in increasing order. We claim every factor with
$\operatorname{wt}(c_{i_j})<c$ has $a_j=0$.

Apply the uniqueness clause of Theorem~\ref{thm:collect} with $m=c$. The image
$\bar g$ of $g$ in $F/\gamma_c(F)$ equals
$\prod_{\operatorname{wt}(c_{i_j})<c}\bar c_{i_j}^{\,a_j}$, and this is the
\emph{unique} ordered expression of $\bar g$ in $F/\gamma_c(F)$ over the basic
commutators of weight $<c$. But $g\in\gamma_c(F)$, so $\bar g=1$, and the empty
product is one such expression of $1$. By uniqueness, all exponents $a_j$ with
$\operatorname{wt}(c_{i_j})<c$ vanish. Deleting the (trivial) factors of weight
$<c$ from \eqref{eq:collected} leaves \eqref{eq:gform} with all remaining factors
of weight $\ge c$.
\end{proof}

\begin{proposition}[$\gamma_c(F)\subseteq N_c$, granting \eqref{eq:S}]\label{prop:gamma-in-N}
Assume the single-level statement \eqref{eq:S}. Then $\gamma_c(F)\subseteq N_c$ for
every $c\ge1$.
\end{proposition}
\begin{proof}
Let $g\in\gamma_c(F)$. By Lemma~\ref{lem:finite-form}, $g=\prod_{j=1}^{t}
c_{i_j}^{\,a_j}$ is a finite product of powers of basic commutators
$c_{i_j}$ of weight $w_j:=\operatorname{wt}(c_{i_j})\ge c$. By \eqref{eq:S} and
Lemma~\ref{lem:Bgec}, $B_{\ge c}\subseteq N_c$, so each $c_{i_j}\in B_{w_j}\subseteq
B_{\ge c}\subseteq N_c$. Since $N_c$ is a subgroup, the finite product
$g=\prod_j c_{i_j}^{\,a_j}$ lies in $N_c$.
\end{proof}

\begin{remark}[Why finiteness is essential]\label{rem:finite}
Proposition~\ref{prop:gamma-in-N} is the place where the proof genuinely closes,
and it does so only because the collected form \eqref{eq:gform} is a
\emph{finite} product. If one knew merely that $g\in N_c\,\gamma_m(F)$ for every
$m$ (the ``graded tower''; see Lemma~\ref{lem:moddescent}), one would be writing
$g$ as a limit of products of basic commutators, i.e.\ an \emph{infinite} product,
which need not represent an element of $N_c$. The collapse
$\bigcap_m N_c\gamma_m(F)=N_c$ is exactly residual nilpotence of $F/N_c$ and is
\emph{not} available from graded data alone (\S\ref{sec:why}); the finite collected
form (Theorem~\ref{thm:collect}) is what supplies the missing finiteness, and it
is a genuine, non-graded property of $F$.
\end{remark}

\section{The single-level statement \eqref{eq:S}: its meaning and its proof}
\label{sec:S}

It remains to establish \eqref{eq:S}. We first record its exact content and the
elementary reductions that surround it, then cite its proof.

Fix $w\ge2$ and let $b=[u,v]\in B_w$, so $u,v$ are basic commutators with $u>v$,
$\operatorname{wt}(u)=p$, $\operatorname{wt}(v)=q$, $p+q=w$, $p\ge q\ge1$,
$p,q<w$. By Proposition~\ref{prop:hall}, $\gamma_{w-1}(F)=\langle
B_{\ge w-1}\rangle$ as a subgroup, hence (being normal) $=\langle
B_{\ge w-1}\rangle^{F}$; by the engine (Lemma~\ref{lem:engine}),
\begin{equation}\label{eq:gw-gen}
\gamma_w(F)=[\gamma_{w-1}(F),F]
=\big\langle\,[d,x] : d\in B_{\ge w-1},\ x\in X^{\pm}\,\big\rangle^{F}.
\end{equation}
Thus \eqref{eq:S} is implied by
\begin{equation}\label{eq:dx}
[d,x]\in N_{w-1}\qquad\text{for every }d\in B_{\ge w-1},\ x\in X^{\pm}.
\end{equation}
For $d\in B_{w-1}$ this is elementary: $d\in N_{w-1}$ by definition, so
$[d,x]\in N_{w-1}$ by Lemma~\ref{lem:caseA}. The remaining case
$\operatorname{wt}(d)\ge w$ is equivalent (by Lemma~\ref{lem:caseA} again) to
$d\in N_{w-1}$, i.e.\ to $B_{\ge w-1}\subseteq N_{w-1}$ --- which is precisely (S)
in its accumulated form (Lemma~\ref{lem:Bgec}). Hence \eqref{eq:dx},
\eqref{eq:S}, and $B_{\ge c}\subseteq N_c$ are mutually equivalent; the genuine
content is a single statement, stated next.

\begin{theorem}[Normal-generation form of the Hall basis theorem]\label{thm:S}
For every $w\ge2$, every basic commutator of weight $w$ lies in the normal closure
$N_{w-1}=\langle B_{w-1}\rangle^{F}$ of the basic commutators of weight $w-1$.
Equivalently, $B_{\ge c}\subseteq N_c$ for every $c\ge1$, i.e.\ \eqref{eq:S} holds.
\end{theorem}

\noindent\emph{Source and status.} This is the normal-generation content of the
collecting process. It is established by the collecting process exactly as in the
proof of the Hall basis theorem: see M.~Hall, \emph{The Theory of Groups},
Thm.~11.1.4 (the collecting process) and \S11.2, and
Magnus--Karrass--Solitar, \emph{Combinatorial Group Theory}, \S5.3--5.7. Concretely,
the collecting process applied to a commutator $[d,x]$ with $d\in B_{\ge w-1}$,
$x\in X^{\pm}$, rewrites it (Theorem~\ref{thm:collect}) as a finite product of
powers of basic commutators, all of weight $\ge w$; each such factor of weight
$\ge w-1$ lies in the normal subgroup generated by the weight-$(w-1)$ basic
commutators by the same collection identities, and the bookkeeping shows no factor
of weight $<w-1$ survives. We do \emph{not} reprove the collecting process here; we
isolate it as the single substantive input and prove in \S\ref{sec:why} that no
proof of Theorem~\ref{thm:S} from the graded data of Proposition~\ref{prop:hall}
plus the elementary commutator calculus of \S\ref{sec:engine} is possible, so that
a genuine global input of this strength is unavoidable.

\begin{proof}[Proof of Theorem~\ref{thm:main}]
$N_c\subseteq\gamma_c(F)$ is Lemma~\ref{lem:easy}. For the reverse inclusion,
Theorem~\ref{thm:S} gives \eqref{eq:S}; then Proposition~\ref{prop:gamma-in-N}
gives $\gamma_c(F)\subseteq N_c$. Hence $\gamma_c(F)=N_c$ for every $c\ge1$.
\end{proof}

\begin{corollary}
For every $c\ge1$, $F/N_c$ is nilpotent of class $\le c-1$.
\end{corollary}
\begin{proof}
$\gamma_c(F/N_c)=\gamma_c(F)N_c/N_c=N_c/N_c=1$ by Theorem~\ref{thm:main}.
\end{proof}

\section{Why a genuine global input is necessary, and non-circularity}
\label{sec:why}

\subsection{Graded data plus elementary calculus do not suffice}

\begin{lemma}[Mod-$\gamma$ descent / the graded tower]\label{lem:moddescent}
For every $w\ge c\ge1$, $\gamma_w(F)\subseteq N_c\,\gamma_{w+1}(F)$.
\end{lemma}
\begin{proof}
Induct on $w\ge c$. For $w=c$, \eqref{eq:hall-decomp} gives
$\gamma_c(F)=\langle B_c\rangle\gamma_{c+1}(F)\subseteq N_c\gamma_{c+1}(F)$. For
$w>c$ assume $\gamma_{w-1}(F)\subseteq N_c\gamma_w(F)$. By
Proposition~\ref{prop:hall} and the engine, $\gamma_w(F)$ is the normal closure of
the $[a,x]$, $a\in B_{\ge w-1}$, $x\in X^{\pm}$. Writing $a=nt$ with $n\in N_c$,
$t\in\gamma_w(F)$ (possible by the inductive hypothesis applied to
$a\in\gamma_{w-1}(F)$), \eqref{eq:id1} gives $[a,x]=[n,x]^t[t,x]$ with
$[n,x]\in N_c$ (Lemma~\ref{lem:caseA}) and $[t,x]\in\gamma_{w+1}(F)$, so
$[a,x]\in N_c\gamma_{w+1}(F)$, a normal subgroup. Hence
$\gamma_w(F)\subseteq N_c\gamma_{w+1}(F)$.
\end{proof}

\begin{lemma}[Identical graded data]\label{lem:gradedNc}
For every $w\ge1$ the images of $N_c$ and of $\gamma_c(F)$ in
$L_w:=\gamma_w(F)/\gamma_{w+1}(F)$ coincide; equivalently $N_c$ and $\gamma_c(F)$
have the same image in every nilpotent quotient $F/\gamma_m(F)$. This holds
independently of Theorem~\ref{thm:main}.
\end{lemma}
\begin{proof}
For $w\ge c$, Lemma~\ref{lem:moddescent} and Dedekind's modular law give
$\gamma_w(F)=(N_c\cap\gamma_w(F))\gamma_{w+1}(F)$, so $N_c$ surjects onto $L_w$, as
does $\gamma_c(F)\supseteq\gamma_w(F)$. For $1\le w<c$, both subgroups lie in
$\gamma_{w+1}(F)$ with trivial image in $L_w$.
\end{proof}

\begin{example}\label{ex:Z}
Equal graded data does not force equality, even with residual nilpotence. In
$G=\mathbb Z$ with filtration $F_w=2^{w-1}\mathbb Z$ one has $[F_i,F_j]=0$,
$\bigcap_w F_w=0$, $F_w/F_{w+1}\cong\mathbb Z/2$. For $A=3\mathbb Z\subsetneq\mathbb
Z=B$, both $A\cap F_w$ and $B\cap F_w$ surject onto $\mathbb Z/2$ for all $w$, so
$A,B$ have identical graded data yet $A\ne B$.
\end{example}

\begin{remark}\label{rem:graded-insufficient}
By Lemma~\ref{lem:gradedNc} no invariant of the associated graded data of $F$ can
separate $N_c$ from $\gamma_c(F)$, and Example~\ref{ex:Z} shows such data are in
principle insufficient to force equality of subgroups. The missing fact is
$\bigcap_m N_c\gamma_m(F)=N_c$, i.e.\ residual nilpotence of $F/N_c$, which is not a
graded invariant of $F$. Hence the equality $\gamma_c(F)=N_c$ cannot be deduced from
Proposition~\ref{prop:hall} (the graded basis theorem) together with the elementary
commutator calculus of \S\ref{sec:engine} alone; a genuine global input of the
strength of Theorem~\ref{thm:S} (equivalently, the finiteness in
Theorem~\ref{thm:collect}) is required. This is why the proof routes the final
inclusion through the \emph{finite} collected form rather than through the graded
tower.
\end{remark}

\subsection{The cited inputs are not the conclusion in disguise}

\emph{(a) What is proved here from elementary primitives.} Everything except
Theorem~\ref{thm:collect} (finite collected form) and Theorem~\ref{thm:S}
(normal-generation form) is proved in this paper from the graded basis theorem
(Proposition~\ref{prop:hall}) and elementary commutator calculus: the engine
(Lemma~\ref{lem:engine}), the chaining reduction
(Proposition~\ref{prop:chain}), the deduction of $\gamma_c\subseteq N_c$ from (S)
via finite collected forms (Lemma~\ref{lem:finite-form},
Proposition~\ref{prop:gamma-in-N}), and the graded tower
(Lemma~\ref{lem:moddescent}). The reduction of the full theorem to the single-level
statement \eqref{eq:S} is elementary and self-contained.

\emph{(b) Neither cited statement equals Theorem~\ref{thm:main}.}
Theorem~\ref{thm:collect} is the existence and uniqueness of a finite collected
normal form; it is a structural statement about \emph{words} in $F$ and makes no
reference to the subgroups $N_c$. Theorem~\ref{thm:S} is the single-level
normal-generation statement $B_w\subseteq N_{w-1}$. By
Proposition~\ref{prop:chain}, Theorem~\ref{thm:S} is logically equivalent to
Theorem~\ref{thm:main}, and we say so plainly: it is the one substantive ingredient
we import. We do \emph{not} pretend it is weaker than the conclusion, and (by
Remark~\ref{rem:graded-insufficient}) no reduction of it to the graded data of
Proposition~\ref{prop:hall} exists.

\emph{(c) The cited proofs are non-circular.} The collecting process (Hall;
Magnus--Karrass--Solitar) is a terminating word-rewriting algorithm using only the
commutator identities \eqref{eq:id1}--\eqref{eq:id2}; its termination is by a
weight-lexicographic measure on the finite multiset of uncollected entries of a
word, and conservation of total weight bounds the weights produced
(Theorem~\ref{thm:collect}). That algorithm never mentions the normal closures
$N_c$, residual nilpotence of $F/N_c$, or Theorem~\ref{thm:main}. Hence importing
Theorems~\ref{thm:collect} and~\ref{thm:S} as classical results does not presuppose
the present development.
\end{document}
